\section{Correctness of Dependency-Closed Durable Acknowledgment}
\label{sec:correctness}

This section formally establishes the correctness of Cstamp-PWAL's
dependency-closed durable acknowledgment condition.  We define the key concepts,
state the safety guarantee, and validate it through deterministic model testing.

\subsection{Definitions}
\label{sec:correctness-definitions}

\subsubsection{Dependency Relation}

We define a read-write dependency relation $\to$ over transactions as follows:

\begin{definition}[Read-Write Dependency]
Transaction $T$ depends on transaction $U$ (written $U \to T$) if:
\begin{enumerate}
  \item $U$ writes to a data item $x$, and
  \item $T$ reads from $x$ after $U$ commits.
\end{enumerate}
\end{definition}

The dependency relation is computed from the read set and write set recorded by
the transaction processing engine at validation time.  Dependencies are directed
edges in a directed acyclic graph (DAG) representing the serialization order of
committed transactions.

\subsubsection{Dependency Frontier}

The dependency frontier of a transaction $T$ is the set of all committed
transactions on which $T$ directly or transitively depends.

\begin{definition}[Dependency Frontier]
The dependency frontier of transaction $T$, denoted $\mathit{DepFrontier}(T)$, is
the smallest set of transactions such that:
\begin{enumerate}
  \item $T \in \mathit{DepFrontier}(T)$, and
  \item For every transaction $U \to T$ in the dependency graph,
        $\mathit{DepFrontier}(U) \subseteq \mathit{DepFrontier}(T)$.
\end{enumerate}
\end{definition}

In Cstamp-PWAL, the dependency frontier is computed as a sound
over-approximation of true transitive dependencies.  That is, it may include
transactions that are not strictly necessary for recovery, but it must not omit
any transaction on which $T$ truly depends.

\subsubsection{Durable Acknowledgment Condition}

\begin{definition}[Dependency-Closed Durable Acknowledgment]
A transaction $T$ is eligible for durable acknowledgment if and only if:
\begin{enumerate}
  \item $T$'s own log record has been persisted to stable storage, and
  \item For every transaction $U \in \mathit{DepFrontier}(T)$ where $U \neq T$,
        $U$'s log record has been persisted to stable storage.
\end{enumerate}
\end{definition}

\subsection{Safety Theorem}
\label{sec:correctness-theorem}

\begin{theorem}[Dependency-Closed Durable Acknowledgment Preserves Recovery Safety]
If a transaction $T$ is acknowledged as durable under the dependency-closed
condition, then a crash at any point after acknowledgment will permit recovery
to replay $T$ together with its entire dependency frontier without data loss or
corruption.  Specifically, if $T$ depends on $U$, then recovery can always find
the state on which $T$ depended.
\end{theorem}

\begin{proof}[Sketch]
By the durable acknowledgment condition, when $T$ is acknowledged, both $T$'s
log record and the log records of all transactions in $\mathit{DepFrontier}(T)$
are persisted.  The dependency frontier is a sound over-approximation: if $T$
truly depends on $U$, then $U \in \mathit{DepFrontier}(T)$, so $U$'s log record
is durable.  At recovery, the WAL layer replays all durable log records in
serialization order (determined by commit timestamps in Cstamp-PWAL).  Since
every transaction that $T$ depends on has its durable log record in the WAL, the
recovery process can reconstruct the state that $T$ observed and apply $T$'s
updates correctly.  Crucially, no acknowledged transaction can be recovered
before transactions it depends on, because we only acknowledge after all
dependencies are durable.
\end{proof}

\subsection{Correctness Validation}
\label{sec:correctness-validation}

We validate the dependency-closed durable acknowledgment condition through
deterministic recovery-safety tests.  The test harness executes a series of
synthetic transactions with controlled read/write patterns, injects crashes at
various points, and verifies that recovery produces the correct state.

Specifically, we test the following scenarios:

\begin{enumerate}
  \item \emph{Sequential dependencies}: Transaction chains where $T_1 \to T_2 \to T_3$.
        We verify that acknowledging $T_3$ does not occur before $T_1$ and $T_2$ are durable.
  \item \emph{Multi-reader scenarios}: Multiple transactions reading from a single
        writer, ensuring that readers are not acknowledged before their writer is durable.
  \item \emph{Sparse dependencies}: Workloads where some transactions are independent.
        We verify that independent transactions can be acknowledged asynchronously
        without waiting for unrelated shards.
  \item \emph{Dependency frontier over-approximation}: Cases where the frontier
        conservatively includes unrelated transactions, ensuring correctness even
        when over-approximation occurs.
\end{enumerate}

The test results confirm that the dependency-closed acknowledgment condition
correctly preserves recovery safety across all tested patterns.

\subsection{Discussion: Over-Approximation and Completeness}
\label{sec:correctness-over-approx}

The dependency frontier is permitted to over-approximate true dependencies.
That is, $\mathit{DepFrontier}(T)$ may include transactions on which $T$ does not
actually depend, as long as it does not \emph{under}-approximate (omit necessary
dependencies).

This property is important for Cstamp-PWAL's design: computing the exact
transitive closure of dependencies on the hot path is expensive.  Instead,
Cstamp-PWAL maintains a conservative over-approximation based on direct
read-write dependencies reported by the transaction processing engine.  On
sparse-dependency workloads, this over-approximation is small and incurs little
extra waiting.  On dense-dependency workloads, it approaches a global-prefix
bound, which is safe though conservative.

